\documentclass[wcp]{jmlr}

\usepackage{longtable}
\usepackage{booktabs}
\usepackage{array}
\usepackage{float}
\usepackage{placeins}
\usepackage{textcomp}
\usepackage[T1]{fontenc}
\usepackage{lmodern}
\usepackage{accsupp}

\DeclareRobustCommand{\Rsq}{\BeginAccSupp{method=hex,unicode,ActualText=005200B2}\ensuremath{R^2}\EndAccSupp{}}

\hypersetup{hidelinks,colorlinks=false,pdfborder={0 0 0}}
\renewcommand{\thetable}{S\arabic{table}}
\renewcommand{\thefigure}{S\arabic{figure}}
\jmlryear{2026}
\jmlrworkshop{ACML 2026}
\jmlrvolume{}
\jmlrpages{}
\makeatletter
\renewcommand*{\@titlefoot}{}
\gdef\@editor{}
\let\cleardoublepage\clearpage
\def\ps@jmlrtps{%
  \let\@mkboth\@gobbletwo
  \def\@oddhead{}%
  \def\@evenhead{}%
  \def\@oddfoot{}%
  \def\@evenfoot{}%
}
\makeatother
\title[Supplementary]{Supplementary Material: Stage-Aware Early Warning of Australian Winter Crop Yield Shortfall}
\author{}

\begin{document}
\maketitle

\section{Fixed-Model and Feature-Regime Details}

\begin{table}[H]
\centering
\caption{Fixed-model lead-time comparison. NYH denotes the no-yield-history weather-soil regime; Op. denotes the operational regime. Values are test RMSE.}
\begin{tabular}{lrrrr}
\toprule
Window & Ridge NYH & Ridge Op. & LGBM NYH & LGBM Op. \\
\midrule
May-Jun & 1.01 & 0.74 & 0.92 & 0.86 \\
May-Jul & 1.00 & 0.72 & 1.06 & 0.85 \\
May-Aug & 1.01 & 0.75 & 0.98 & 0.89 \\
May-Sep & 0.96 & 0.70 & 0.94 & 0.89 \\
May-Oct & 0.89 & 0.66 & 0.91 & 0.83 \\
\bottomrule
\end{tabular}

\end{table}

\begin{table}[H]
\centering
\caption{Detailed validation-selected lead-time performance by forecast window for the no-yield-history weather-soil and operational regimes. This table provides the numeric rows summarized by the main-paper lead-time RMSE comparison figure.}
\resizebox{\textwidth}{!}{\begin{tabular}{@{}lllrrr@{}}
\toprule
Window & Feature set & Model & MAE & RMSE & \Rsq{} \\
\midrule
May-Jun & NYH weather-soil & XGBoost & 0.651 & 0.914 & 0.478 \\
May-Jul & NYH weather-soil & CatBoost & 0.659 & 0.951 & 0.436 \\
May-Aug & NYH weather-soil & HistGradientBoosting & 0.609 & 0.968 & 0.416 \\
May-Sep & NYH weather-soil & SVR-RBF & 0.609 & 0.881 & 0.516 \\
May-Oct & NYH weather-soil & SVR-RBF & 0.548 & 0.821 & 0.579 \\
May-Jun & Operational & ElasticNet & 0.533 & 0.719 & 0.677 \\
May-Jul & Operational & Ridge & 0.527 & 0.716 & 0.680 \\
May-Aug & Operational & ElasticNet & 0.506 & 0.724 & 0.673 \\
May-Sep & Operational & Ridge & 0.469 & 0.704 & 0.691 \\
May-Oct & Operational & Ridge & 0.428 & 0.660 & 0.728 \\
\bottomrule
\end{tabular}
}
\end{table}

\begin{table}[H]
\centering
\caption{Feature groups used for stage-aware early-warning models.}
\resizebox{\textwidth}{!}{\begin{tabular}{lll}
\toprule
Feature group & Examples & Interpretation \\
\midrule
Rainfall/dryness & rain\_sum, rain\_days, dry\_days, dry-spell length & Water supply and dry-spell stress \\
Heat/cold & tmax\_mean, heat\_days\_30, heat degree days, frost days & Heat exposure and frost/cold risk \\
Energy/evaporation & radiation\_sum, evap\_sum, vapor pressure & Growth energy and evaporative demand \\
Compound stress & hot-dry days, high-evaporation dry days & Combined weather stress \\
Soil background & AWC, clay/sand, SOC, pH, bulk density, depth & Regional vulnerability conditioning \\
\bottomrule
\end{tabular}
}
\end{table}

\begin{table}[H]
\centering
\caption{Feature-regime definitions used for ablation and information-source analysis.}
\resizebox{\textwidth}{!}{\begin{tabular}{lllll}
\toprule
Feature set & Weather & Soil & Yield hist. & Interpretation \\
\midrule
Identity/time only & No & No & No & Crop-region persistence benchmark \\
Weather & Yes & No & No & Stage-weather signal only \\
Weather+Dev & Yes & No & No & Adds train-derived anomaly features \\
Weather+Soil & Yes & Yes & No & Adds static regional vulnerability \\
No-yield-history weather-soil & Yes & Yes & No & Main information-isolation regime \\
Yield history & No & No & Yes & Persistent productivity memory \\
Operational with yield history & Yes & Yes & Yes & Main deployment-oriented regime \\
\bottomrule
\end{tabular}
}
\end{table}

\begin{figure}[H]
\centering
\includegraphics[width=\textwidth]{figures/fig13_ablation_r2_delta_by_feature_set.png}
\caption{Ablation heatmap showing \Rsq{} differences relative to the full operational feature set. Short labels are used to keep the dense comparison readable.}
\end{figure}
\FloatBarrier

\section{Data Coverage and Baseline Suite}

\begin{table}[H]
\centering
\caption{Missing source yield entries by state-crop series. The total confirms 966 observed entries out of 990 possible state-crop-year combinations in 1989--2021 for the observed crop-region pairs.}
\begin{tabular}{llrr}
\toprule
Region & Crop & Observed years & Missing years \\
\midrule
New South Wales & Barley & 33 & 0 \\
New South Wales & Canola & 33 & 0 \\
New South Wales & Lupins & 33 & 0 \\
New South Wales & Oats & 33 & 0 \\
New South Wales & Wheat & 33 & 0 \\
Queensland & Barley & 33 & 0 \\
Queensland & Canola & 30 & 3 \\
Queensland & Lupins & 24 & 9 \\
Queensland & Oats & 33 & 0 \\
Queensland & Wheat & 33 & 0 \\
South Australia & Barley & 33 & 0 \\
South Australia & Canola & 33 & 0 \\
South Australia & Lupins & 33 & 0 \\
South Australia & Oats & 33 & 0 \\
South Australia & Wheat & 33 & 0 \\
Tasmania & Barley & 33 & 0 \\
Tasmania & Canola & 33 & 0 \\
Tasmania & Lupins & 21 & 12 \\
Tasmania & Oats & 33 & 0 \\
Tasmania & Wheat & 33 & 0 \\
Victoria & Barley & 33 & 0 \\
Victoria & Canola & 33 & 0 \\
Victoria & Lupins & 33 & 0 \\
Victoria & Oats & 33 & 0 \\
Victoria & Wheat & 33 & 0 \\
Western Australia & Barley & 33 & 0 \\
Western Australia & Canola & 33 & 0 \\
Western Australia & Lupins & 33 & 0 \\
Western Australia & Oats & 33 & 0 \\
Western Australia & Wheat & 33 & 0 \\
Total & All crops & 966 & 24 \\
\bottomrule
\end{tabular}

\end{table}

\tiny
\begin{longtable}{@{}p{0.15\textwidth}p{0.21\textwidth}p{0.20\textwidth}p{0.08\textwidth}rrr@{}}
\caption{Full May-Oct strong-baseline suite under the same leakage-safe split. External crop-yield studies are used as methodological context, not as numerical baselines, because their crops, spatial scales, inputs, lead times, and evaluation protocols differ.}\label{tab:baselines-full-appendix}\\
\toprule
Group & Regime & Model & Yield hist. & MAE & RMSE & \Rsq{} \\
\midrule
\endfirsthead
\toprule
Group & Regime & Model & Yield hist. & MAE & RMSE & \Rsq{} \\
\midrule
\endhead
Historical & Past-only yield & 3-year rolling mean & Yes & 0.606 & 0.779 & 0.622 \\
Historical & Past-only yield & Previous-year yield & Yes & 0.562 & 0.784 & 0.616 \\
Historical & Past-only yield & Crop-region train trend & Yes & 0.613 & 0.823 & 0.577 \\
Historical & Past-only yield & Crop-region train mean & Yes & 0.683 & 0.961 & 0.423 \\
Classical ML & NYH weather-soil & SVR-RBF & No & 0.548 & 0.821 & 0.579 \\
Classical ML & NYH weather-soil & Ridge & No & 0.536 & 0.889 & 0.507 \\
Classical ML & NYH weather-soil & ElasticNet & No & 0.557 & 0.906 & 0.487 \\
Classical ML & NYH weather-soil & RandomForest & No & 0.685 & 1.012 & 0.360 \\
Classical ML & Operational & Ridge & Yes & 0.428 & 0.660 & 0.728 \\
Classical ML & Operational & ElasticNet & Yes & 0.439 & 0.671 & 0.719 \\
Classical ML & Operational & SVR-RBF & Yes & 0.500 & 0.757 & 0.642 \\
Classical ML & Operational & RandomForest & Yes & 0.564 & 0.836 & 0.564 \\
Strong tabular ML & NYH weather-soil & XGBoost & No & 0.583 & 0.888 & 0.508 \\
Strong tabular ML & NYH weather-soil & CatBoost & No & 0.591 & 0.901 & 0.493 \\
Strong tabular ML & NYH weather-soil & LightGBM & No & 0.581 & 0.907 & 0.486 \\
Strong tabular ML & NYH weather-soil & HistGradientBoosting & No & 0.605 & 0.925 & 0.466 \\
Strong tabular ML & Operational & CatBoost & Yes & 0.505 & 0.778 & 0.622 \\
Strong tabular ML & Operational & XGBoost & Yes & 0.515 & 0.779 & 0.621 \\
Strong tabular ML & Operational & HistGradientBoosting & Yes & 0.521 & 0.812 & 0.588 \\
Strong tabular ML & Operational & LightGBM & Yes & 0.532 & 0.826 & 0.575 \\
Interpretable ML & NYH weather-soil & GAM & No & 0.675 & 1.045 & 0.318 \\
Interpretable ML & Operational & GAM & Yes & 0.646 & 1.159 & 0.162 \\
Sequence & Daily sequence NYH & DailyWeather-GRU & No & 0.674 & 1.020 & 0.350 \\
Sequence & Daily sequence operational & DailyWeather-GRU & Yes & 0.520 & 0.718 & 0.678 \\
\bottomrule
\end{longtable}
\normalsize


All-window baseline-suite metrics remain available in the accompanying CSV artifacts; the PDF reports the full May-Oct suite to keep the comparison readable.

\begin{table}[H]
\centering
\caption{Detailed stress-validation summary by validation protocol and feature regime. Rows report mean performance across the folds and forecast windows belonging to each protocol; the held-out-test row is an all-window stress summary, not the May-Oct-only result in the main lead-time comparison.}
\resizebox{\textwidth}{!}{\begin{tabular}{lllrrr}
\toprule
Protocol & Feature set & Model & Mean RMSE & Mean \Rsq{} & Folds \\
\midrule
Leave-one-crop-out & Operational with yield history & Ridge & 0.731 & 0.520 & 5 \\
Leave-one-crop-out & No-yield-history weather-soil & Ridge & 1.071 & 0.067 & 5 \\
Leave-one-region-out & Operational with yield history & Ridge & 0.823 & -0.356 & 6 \\
Leave-one-region-out & No-yield-history weather-soil & LightGBM & 0.974 & -0.285 & 6 \\
Rolling-origin & Operational with yield history & Ridge & 0.516 & 0.698 & 3 \\
Rolling-origin & No-yield-history weather-soil & LightGBM & 0.578 & 0.631 & 3 \\
Held-out test years, all-window mean & Operational with yield history & Ridge & 0.716 & 0.679 & 1 \\
Held-out test years, all-window mean & No-yield-history weather-soil & LightGBM & 0.963 & 0.420 & 1 \\
\bottomrule
\end{tabular}
}
\end{table}

\section{Data and Implementation Details}

The raw yield source is the ABARES Australian Crop Report state-crop-data workbook \path{03_AustCropRpt20260303_StateCropData_v1.0.0.xlsx}, accessed on 4 July 2026; the modeled panel is the derived \texttt{yield\_panel.csv}. Fields used are crop area, production, and yield in tonnes per hectare, filtered to the five winter crops and six state-level regions in the main paper. The harvest year is represented by \texttt{year\_start}. Production and area are retained for audit but excluded from model features.

SILO/LongPaddock weather inputs are daily rainfall, maximum and minimum temperature, radiation, evaporation, and vapor pressure. They are aggregated from May through each forecast cutoff. The system uses state-level regional daily series as a monitoring exposure proxy; latitude and longitude metadata are retained, but no crop-area-weighted gridded exposure is claimed. Weather-deviation features use training years only: for feature value $x_{r,w,t}$, the deviation is $x_{r,w,t}-\bar{x}_{r,w}^{train}$ for the same region and forecast window. Soil/ASRIS attributes are static regional background covariates from the Soil and Landscape Grid of Australia, including available water capacity, texture, bulk density, soil organic carbon, nutrients, pH, cation exchange capacity, and depth variables. These features condition regional vulnerability and may partly proxy region identity.

The no-yield-history regime has 73 predictors before one-hot expansion: crop, state-level region, year, 30 stage-weather summaries, 18 train-derived weather-deviation features, and 22 soil background attributes after excluding coordinate metadata. The operational regime adds one-year lag, three-year rolling past mean, and expanding past mean, for 76 predictors before one-hot expansion. Scaling, imputation, encoders, baselines, thresholds, and conformal residuals are fitted without test years; one-hot encoders handle unknown categories in stress tests. Configured weather thresholds are rain day 1 mm, heavy rain 10/25 mm, heat 25/30/35 C, frost 0 C, cold 5 C, and high evaporation 5 mm.

\section{Baseline-Suite Implementation}

The internal baseline suite is run once with random seed 42. This single-seed design is a limitation; model selection is by validation RMSE within each forecast window and feature regime, and test metrics are reported once. Tabular preprocessing uses median imputation and standard scaling for numeric predictors, one-hot encoding for crop and region, and unknown-category handling in held-out stress tests. The suite uses scikit-learn 1.9.0, XGBoost 3.3.0, LightGBM 4.6.0, CatBoost 1.2.10, pyGAM 0.12.0, and PyTorch 2.12.1+cpu in the recorded run.

Candidate models include Ridge with alpha 1.0 or 3.0; ElasticNet with alpha 0.01/l1 ratio 0.2 or alpha 0.03/l1 ratio 0.4; RandomForest with 400 trees and either unrestricted depth/min leaf 3 or depth 8/min leaf 5; SVR-RBF with C 1 or 10 and epsilon 0.1; HistGradientBoosting; XGBoost; LightGBM; CatBoost; and a pyGAM additive model over transformed features. The GRU sequence comparator uses daily weather up to the cutoff window, static covariates concatenated after the sequence encoder, and for the operational regime adds lag/rolling/expanding yield-history features to the static block. Two GRU configurations are validation-selected: hidden size 16 or 32, dropout 0.15 or 0.20, Adam learning rate 0.001, weight decay 1e-4, up to 160 epochs with patience 16.

\section{Classification and Uncertainty Details}

\begin{table}[H]
\centering
\caption{Sensitivity of the low-yield-risk label to crop-specific training shortfall percentiles. The 80th percentile is the main setting.}
\begin{tabular}{rrrr}
\toprule
Percentile & Train rate (\%) & Validation rate (\%) & Test rate (\%) \\
\midrule
70 & 30.1 & 13.8 & 24.2 \\
80 & 20.3 & 11.2 & 18.1 \\
90 & 10.4 & 2.6 & 9.4 \\
\bottomrule
\end{tabular}

\end{table}

\begin{table}[H]
\centering
\caption{Best threshold-tuned classification metrics by forecast window. Pred.+ is the fraction of test examples flagged positive by the selected threshold.}
\resizebox{\textwidth}{!}{\begin{tabular}{lllrrrrrrrr}
\toprule
Window & Model & Strategy & Thr. & Prec. & Recall & F1 & ROC-AUC & PR-AUC & Brier & Pred.+ \\
\midrule
May-Jun & CatBoost & best F1 & 0.14 & 0.21 & 0.89 & 0.33 & 0.51 & 0.19 & 0.149 & 0.79 \\
May-Jul & LightGBM & best F1 & 0.30 & 0.27 & 0.70 & 0.39 & 0.63 & 0.25 & 0.144 & 0.47 \\
May-Aug & LogisticRegression & precision $\geq$ 0.3 & 0.22 & 0.27 & 0.78 & 0.40 & 0.65 & 0.25 & 0.147 & 0.52 \\
May-Sep & HistGradientBoosting & best F1 & 0.63 & 0.28 & 0.74 & 0.41 & 0.69 & 0.27 & 0.197 & 0.48 \\
May-Oct & LogisticRegression & precision $\geq$ 0.3 & 0.27 & 0.47 & 0.56 & 0.51 & 0.80 & 0.39 & 0.135 & 0.21 \\
\bottomrule
\end{tabular}
}
\end{table}

\begin{table}[H]
\centering
\caption{Low-yield-risk confusion matrix using validation-selected thresholds.}
\resizebox{\textwidth}{!}{\begin{tabular}{lllrrrrrrr}
\toprule
Window & Model & Thr. & TN & FP & FN & TP & Prec. & Recall & F1 \\
\midrule
May-Oct & LogisticRegression & 0.27 & 105 & 17 & 12 & 15 & 0.469 & 0.556 & 0.508 \\
\bottomrule
\end{tabular}
}
\end{table}

\begin{table}[H]
\centering
\caption{Held-out May-Oct ranked watch-list example. Rank scores are uncalibrated classifier scores used for ordering analyst-review lists and should not be interpreted as fully calibrated event probabilities. Observed indicates the held-out low-yield-risk label.}
\resizebox{\textwidth}{!}{\begin{tabular}{rllrrrrr}
\toprule
Rank & Region & Crop & Year & Rank score & Alert & Observed & Shortfall \\
\midrule
1 & VIC & Oats & 2018 & 1.000 & 1 & 1 & 0.465 \\
2 & NSW & Lupins & 2018 & 1.000 & 1 & 1 & 0.410 \\
3 & VIC & Lupins & 2018 & 1.000 & 1 & 0 & 0.274 \\
4 & VIC & Canola & 2018 & 1.000 & 1 & 0 & 0.012 \\
5 & NSW & Wheat & 2018 & 0.273 & 1 & 1 & 0.979 \\
6 & QLD & Wheat & 2019 & 0.273 & 1 & 1 & 0.957 \\
7 & QLD & Wheat & 2018 & 0.273 & 1 & 1 & 0.880 \\
8 & QLD & Barley & 2019 & 0.273 & 1 & 1 & 0.865 \\
9 & TAS & Oats & 2019 & 0.273 & 1 & 1 & 0.730 \\
10 & QLD & Canola & 2018 & 0.273 & 1 & 1 & 0.646 \\
\bottomrule
\end{tabular}
}
\end{table}

Coverage below the nominal 80\% target, especially in later windows, is interpreted as temporal calibration difficulty under small validation samples. These intervals are diagnostics for analyst review rather than reliable calibrated intervals for automatic decisions.

\begin{table}[H]
\centering
\caption{Compact conformal uncertainty diagnostics by forecast window. Coverage and interval width are evaluated on held-out test years after validation-period calibration.}
\resizebox{\textwidth}{!}{\begin{tabular}{llrrrrr}
\toprule
Window & Interval model & Nominal & Coverage & Width & qhat & Brier \\
\midrule
May-Jun & Conformal ElasticNet & 0.800 & 0.799 & 1.565 & 0.782 & 0.149 \\
May-Jul & Conformal ElasticNet & 0.800 & 0.792 & 1.423 & 0.711 & 0.144 \\
May-Aug & Conformal ElasticNet & 0.800 & 0.711 & 1.259 & 0.629 & 0.147 \\
May-Sep & Conformal ElasticNet & 0.800 & 0.758 & 1.251 & 0.626 & 0.197 \\
May-Oct & Conformal ElasticNet & 0.800 & 0.752 & 1.125 & 0.563 & 0.135 \\
\bottomrule
\end{tabular}
}
\end{table}

\section{Supplementary Figures}

\begin{figure}[H]
\centering
\includegraphics[width=0.72\textwidth]{figures/fig14_stress_validation_heatmap.png}
\caption{Stress validation heatmap by validation protocol and feature regime. Values use the same best-model, all-window summary as the main-paper stress validation table, not single-window May-Oct scores.}
\end{figure}

\begin{figure}[H]
\centering
\includegraphics[width=0.94\textwidth]{figures/fig19_feature_importance_split.png}
\caption{Feature group importance measured as mean RMSE increase after group permutation. Operational importance should not be interpreted as within-season weather importance because lag-yield features are allowed in that regime.}
\end{figure}

\end{document}
